\documentclass[10pt,a4paper]{amsart}
\usepackage[margin=19mm]{geometry}
\usepackage{amsmath,amssymb,mathtools}
\usepackage[scr=rsfs,cal=euler]{mathalfa}
\usepackage{xfrac}
\usepackage[unicode,hidelinks,bookmarks=false]{hyperref}
\newcommand{\ie}{i.e.\ }
\newcommand{\cf}{cf.\ }
\newcommand{\resp}{respectively\ }
\newcommand{\Subsection}{\subsection}
\providecommand{\nobreakdash}{\nobreak\hbox{-}}
\setlength{\emergencystretch}{2em}
\multlinegap0pt
\usepackage{amssymb}
\usepackage{algorithm}\usepackage{algpseudocode}
\usepackage{xcolor}
\usepackage[shortlabels]{enumitem}
\newtheorem{lemma}{Lemma}
\newtheorem{proposition}{Proposition}
\newtheorem{theorem}{Theorem}
\title{Beyond Polynomials: Optimal Locally Recoverable Codes from Good Rational Functions}
\author{Hengfeng Liu, Sihem Mesnager, Chunming Tang, Xuemin Zheng}
\date{Source: arXiv:2605.11465v1 (CC BY 4.0)}
\begin{document}
\maketitle

\noindent\emph{Source and licence.} Beyond Polynomials: Optimal Locally Recoverable Codes from Good Rational Functions. Version arXiv:2605.11465v1 (CC BY 4.0), distributed by arXiv under the Creative Commons Attribution 4.0 International licence, http://creativecommons.org/licenses/by/4.0/, as recorded in the arXiv licence metadata for this exact version and shown on the abstract page https://arxiv.org/abs/2605.11465v1. Full licence notice: \texttt{licenses/arxiv-2605.11465-CC-BY-4.0.txt}.\par\smallskip

\noindent\textit{Editorial scope.} Counted unit: the article's theorem 21e (source0.tex lines 991--1004). The article's complete original proof of it is delivered with it (source0.tex lines 1006--1105, one \texttt{proof} environment of 100 lines). No mathematics is written, repaired or re-derived: the statement and the proof are the article's own lines, in the article's own notation, and no retained line is substituted or re-flowed.  Retained uncounted as the source-local context the proof needs: lines 680--686; lines 976--985.  Nothing was omitted from any retained span, so the package carries no omission record.  No retained line contains a citation, so the wrapper carries no bibliography and no bibliography entry.  No retained line contains a figure environment, an image-inclusion command or drawing code, so no image is delivered and the package declares no asset dependency.  Editorial formatting only, in addition: an AMS \texttt{amsart} preamble that restates the article's own declarations and macros used by the retained lines, namely the environments \texttt{lemma}, \texttt{proposition}, \texttt{theorem} on separate counters, together with the standard packages those lines need.  The retained environments print in this document as Lemma 1, Proposition 1, Theorem 1 in order of appearance, whereas the article prints the same items with the numbering of its own full document.  The complete native source of the article as distributed in the arXiv e-print for this exact version is bundled unchanged as \texttt{sources/200360/source0.tex}.
\begin{lemma}\label{M}
Let $h(x)$ be a rational function, and let $t=h(x)$. Let $M_{h}$ be the Galois closure of the extension $\mathbb{F}_q(x)/\mathbb{F}_q(t)$. Then for any place $P$ in $\mathbb{F}_{q}(t)$, we have
\begin{itemize}
\item[(i)] $P$ splits totally in the extension $\mathbb{F}_q(x)/\mathbb{F}_q(t)$ if and only if it splits totally in the extension $M_{h}/\mathbb{F}_q(t)$. 
\item[(ii)] $P$ is ramified in the extension $\mathbb{F}_q(x)/\mathbb{F}_q(t)$ if and only if it is ramified in the extension $M_{h}/\mathbb{F}_q(t)$. 
\end{itemize}
\end{lemma}

\begin{proposition}\label{infi}
Let $h(x)=\displaystyle\frac{f(x)}{g(x)}$ be a rational function with $f(x)$ and $g(x)$ coprime. Let the factorization of $g(x)$ in $\mathbb{F}_{q}\left [ X \right ] $ be $g(x)=\displaystyle\prod_{ i }g_{{i}}(x)^{a_{i}}$, where the $g_{i}(x)$ are irreducible. Denote by $P_{g_{i}}$ the place associated with $g_{i}(x)$ in $\mathbb{F}_{q}(x)$, and let $P_{\infty}^{t}$ (resp., $P_{\infty}^{x}$ ) be the place at infinity in $\mathbb{F}_{q}(t)$ (resp., $\mathbb{F}_{q}(x)$). 

Let $h(x)=t$. In the extension $\mathbb{F}_{q}(x)/\mathbb{F}_{q}(t)$, the places of $\mathbb{F}_{q}(x)$ lying over $P_{\infty}^{t}$, together with their ramification indices and relative degrees, are given as follows:

\begin{itemize}
\item[(i)] If $\deg(f)\le \deg(g)$, then the places lying over $P_{\infty}^{t}$ are the $P_{g_{i}}$, with $e(P_{g_{i}}|P_{\infty}^{t})=a_{i}$ and $f(P_{g_{i}}|P_{\infty}^{t})=\deg(g_{i})$.
\item[(ii)] If $\deg(f)>\deg(g)$, then the places lying over $P_{\infty}^{t}$ are the $P_{g_{i}}$ and $P_{\infty}^{x}$, with $e(P_{g_{i}}|P_{\infty}^{t})=a_{i}$, $f(P_{g_{i}}|P_{\infty}^{t})=\deg(g_{i})$, $e(P_{\infty}^{x}|P_{\infty}^{t})=\deg(f)-\deg(g)$, and $f(P_{\infty}^{x}|P_{\infty}^{t})=1$. 
\end{itemize} 
\end{proposition}

\begin{theorem} \label{21e}
Let $h(x)=\displaystyle\frac{f(x)}{g(x)}$ be a rational function with $f(x)$ and $g(x)$ coprime. Let $t=h(x)$, and let $M_{h}$ be the Galois closure of the extension $\mathbb{F}_q(x)/\mathbb{F}_q(t)$. In the extension $M_{h}/\mathbb{F}_q(t)$, let $R_{h}$ (resp., $R_{h}^{1}$) be the set of all ramified places (resp., ramified rational places). Let $g_{h}$ be the genus of $M_{h}$, and let $G_{h}$ be the Galois group of $M_{h}/\mathbb{F}_q(t)$. Then 

\begin{itemize}
\item[(i)] $\#R_{h}^{1}\le \#R_{h}\le \deg(h)+\displaystyle\sum_{i}\deg(g_{i})+\delta_{g}^{f}-1$, where $\deg(h)=\max\{deg(f), deg(g)  \}$, and
$$\delta_{g}^{f}=\begin{cases}
1 & \deg(f)>\deg(g), \\
0 & \deg(f)\le\deg(g).
\end{cases}$$

\item[(ii)] Suppose $\text{char}(\mathbb{F}_{q})\nmid \#G_{h}$ and the constant field of $M_{h}$ is $\mathbb{F}_{q}$. Then
$$g_{h}\le (\deg(h)-2)\#G_{h}+1.$$
\end{itemize}
\end{theorem}

\begin{proof}
We begin with the proof of $(i)$.

\medskip

By Lemma \ref{M}, the number $\#R_{h}$ coincides with the number of ramified places in the extension $\mathbb{F}_q(x)/\mathbb{F}_q(t)$. Let $\mathbb{P}(\mathbb{F}_{q})$ denote the set of all places of $\mathbb{F}_{q}(t)$. Applying the Riemann--Hurwitz formula to the extension $\mathbb{F}_q(x)/\mathbb{F}_q(t)$, we obtain
\begin{equation}\label{H1}
\begin{aligned}
2\deg(h)-2
\ge
\sum_{P\in \mathbb{P}(\mathbb{F}_{q}) }\sum_{P'|P}(e(P'|P)-1)\deg(P')
=
\sum_{P\in \mathbb{P}(\mathbb{F}_{q}) }\delta_{P}.
\end{aligned}
\end{equation}

Separating the contribution of the infinite place, we write
\[
\sum_{P\in \mathbb{P}(\mathbb{F}_{q}) }\delta_{P}
=
\sum_{\substack{P\ne P_{\infty} \\ \text{$P$ ramified}}} \delta_{P}
+
\delta_{\infty}.
\]

Using Proposition \ref{infi}, we compute
\begin{equation}\label{H2}
\begin{aligned}
\delta_{\infty}
&=
\sum_{P'|P_{\infty}}(e(P'|P_{\infty})-1)f(P'|P_{\infty}) \\
&=
\sum_{P'|P_{\infty}}e(P'|P_{\infty})f(P'|P_{\infty})
-
\sum_{P'|P_{\infty}}f(P'|P_{\infty}) \\
&=
\deg(h)-\sum_{i}\deg(g_{i})-\delta_{g}^{f}.
\end{aligned}
\end{equation}

Combining (\ref{H1}) and (\ref{H2}), we obtain
\[
\#R_{h}
\le
\sum_{\substack{P\ne P_{\infty} \\ \text{$P$ ramified}}} \delta_{P}+1
\le
2\deg(h)-\delta_{\infty}-1
=
\deg(h)+\sum_{i}\deg(g_{i})+\delta_{g}^{f}-1.
\]

This proves $(i)$.

\medskip

We now prove $(ii)$.

\medskip

We apply the Riemann--Hurwitz formula to the Galois extension $M_{h}/\mathbb{F}_q(t)$. Since $\text{char}(\mathbb{F}_{q})\nmid \#G_{h}$ and the constant field of $M_{h}$ is $\mathbb{F}_{q}$, the extension is tame, and equality holds. Hence,
\begin{equation}\label{ii1}
\begin{aligned}
2g_{h}-2
&=
-2\#G_{h}
+
\sum_{P\in \mathbb{P}(\mathbb{F}_{q}) }\deg(P)\sum_{P'|P}(e(P'|P)-1)f(P'|P).
\end{aligned}
\end{equation}

Using the transitivity of the Galois action, we obtain
\[
\sum_{P'|P}(e(P'|P)-1)f(P'|P)
=
(e(P)-1)f(P)\cdot \frac{\#G_{h}}{e(P)f(P)}.
\]

Substituting into (\ref{ii1}), we deduce
\[
2g_{h}-2
\le
\#G_{h}\left ( \sum_{P\in \mathbb{P}(\mathbb{F}_{q}) }\deg(P)-2 \right ).
\]

To bound $\sum_{P}\deg(P)$, we use again (\ref{H1}), which yields
\begin{equation}\label{ii2}
\sum_{P\in \mathbb{P}(\mathbb{F}_{q}) }\deg(P)
\le
\sum_{P\in \mathbb{P}(\mathbb{F}_{q}) }\delta_{P}
\le
2\deg(h)-2.
\end{equation}

Combining (\ref{ii1}) and (\ref{ii2}), we conclude that
\[
g_{h}\le (\deg(h)-2)\#G_{h}+1.
\]

This completes the proof.
\end{proof}

\end{document}
